\sdocumentclass[11pt]{article}
\usepackage{times}
\usepackage[left=1in,top=1in,right=1in,bottom=1in,nohead,nofoot]{geometry}

\begin{document}

\thispagestyle{empty}
%\section*{A. Project Summary}
%\section*{Project Summary}
%\vspace{0.2in}

\section*{Overview}

% Cache and memory systems are increasingly complex, and thus ML techniques
% can play increasingly important roles

Cache and main memory are critical to the performance and energy efficiency of
today's computer systems.  With their increasing design complexity, it is more
and more challenging to optimize their performance for memory-intensive
workloads.  Machine learning (ML) techniques have been explored in a limited
scope to improve their performance, particularly in cache prefetching and memory
access scheduling. However, there remains considerable opportunities for ML
techniques to play a more substantial role.

% Here discuss the scope of ML components
This project will investigate a holistic framework of ML techniques to improve
the performance and energy efficiency of caches and main memories for multi-core
processors.  It consists of three ML components for cache replacement, cache
prefetch filtering, and memory access scheduling, respectively.  Those ML
components will be seamlessly integrated with the corresponding hardware
components, namely cache controller, cache prefetcher, and memory controller, to
refine their decision-making logic for better performance and energy efficiency.
Those ML components will improve cache hit rates, prefetching accuracy, and
memory channel utilization, respectively.  Working together, they may
significantly improve the performance and energy efficiency of a wide range of
memory intensive workloads, without any significant increase the overall cost
of the processors.

% Here discuss the design of each component
The project will study in depth the design of those ML components, focusing on
circuit efficiency and ML robustness.  Instead of using static feature selection
and weight tables customized for pre-selected and limited features, in this
project the ML components will each use a set of shared weight tables with a
much larger set of selectable features.  Algorithms will be developed to choose
an optimal set of features dynamically for the current workloads and their
current running phases. This approach enables machine learning over a broad
feature set without excessive hardware cost, and therefore can be explore more
opportunities for machine learning.  The project will investigate the design and
optimization space starting from a baseline design, and then evaluate the ML
learning effectiveness of those components as well as the overall improvements
of performance and energy efficiency.

% Here discuss the coordination
The project will further develop a coordination framework for those ML
components, which are distributed within a multi-core processor, to improve
their learning effectiveness.  This research direction explores potential,
inherent correlations among caching, prefetching, and memory access scheduling
to coordinate the learning process of those distributed ML components.  In the
coordination framework, they may interactively learn from each other about
workload phase changes, optimal dynamic feature selection, optimal learning
intensity, and others; and may provide new features to broaden the feature
selection.  The project will fully investigate the potential of this approach
and its limitations.

\section*{Intellectual Merits}

The proposed research presents innovative strategies of applying machine
learning techniques to improve the performance and energy efficiency of cache
and main memory systems.  It incorporates distributed ML components throughout
the cache and memory hierarchy, and coordinates their learning processes
systematically.  It uses dynamic feature selection to broaden the feature set
without excessive hardware cost, and allows those ML components to learn from
each other in coordination.  The approach will provide new insights about how to
utilize ML techinques effectively in improving memory performance and energy
efficiency.

\section*{Broader Impacts}

The proposed research may significantly improve the performance and energy
efficiency of modern cache and memory systems by leveraging the recent advances
in machine learning techniques.  The education components will introduce the
timely topics of memory system design and memory system reliability into
undergraduate and graduate courses, broadening the knowledge of future computer
engineers.

\end{document}
